% Options for packages loaded elsewhere
\PassOptionsToPackage{unicode}{hyperref}
\PassOptionsToPackage{hyphens}{url}
%
\documentclass[
]{article}
\usepackage{amsmath,amssymb}
\usepackage{lmodern}
\usepackage{iftex}
\ifPDFTeX
  \usepackage[T1]{fontenc}
  \usepackage[utf8]{inputenc}
  \usepackage{textcomp} % provide euro and other symbols
\else % if luatex or xetex
  \usepackage{unicode-math}
  \defaultfontfeatures{Scale=MatchLowercase}
  \defaultfontfeatures[\rmfamily]{Ligatures=TeX,Scale=1}
\fi
% Use upquote if available, for straight quotes in verbatim environments
\IfFileExists{upquote.sty}{\usepackage{upquote}}{}
\IfFileExists{microtype.sty}{% use microtype if available
  \usepackage[]{microtype}
  \UseMicrotypeSet[protrusion]{basicmath} % disable protrusion for tt fonts
}{}
\makeatletter
\@ifundefined{KOMAClassName}{% if non-KOMA class
  \IfFileExists{parskip.sty}{%
    \usepackage{parskip}
  }{% else
    \setlength{\parindent}{0pt}
    \setlength{\parskip}{6pt plus 2pt minus 1pt}}
}{% if KOMA class
  \KOMAoptions{parskip=half}}
\makeatother
\usepackage{xcolor}
\usepackage{longtable,booktabs,array}
\usepackage{multirow}
\usepackage{calc} % for calculating minipage widths
% Correct order of tables after \paragraph or \subparagraph
\usepackage{etoolbox}
\makeatletter
\patchcmd\longtable{\par}{\if@noskipsec\mbox{}\fi\par}{}{}
\makeatother
% Allow footnotes in longtable head/foot
\IfFileExists{footnotehyper.sty}{\usepackage{footnotehyper}}{\usepackage{footnote}}
\makesavenoteenv{longtable}
\setlength{\emergencystretch}{3em} % prevent overfull lines
\providecommand{\tightlist}{%
  \setlength{\itemsep}{0pt}\setlength{\parskip}{0pt}}
\setcounter{secnumdepth}{-\maxdimen} % remove section numbering
\ifLuaTeX
  \usepackage{selnolig}  % disable illegal ligatures
\fi
\IfFileExists{bookmark.sty}{\usepackage{bookmark}}{\usepackage{hyperref}}
\IfFileExists{xurl.sty}{\usepackage{xurl}}{} % add URL line breaks if available
\urlstyle{same} % disable monospaced font for URLs
\hypersetup{
  hidelinks,
  pdfcreator={LaTeX via pandoc}}

\author{}
\date{}

\begin{document}

\textbf{An Automated Pipeline for Flagging Outdated Citations in
AI-Recommended Literature: A Comprehensive Review}

\begin{quote}
RITAM SANTRA, Indian Institute of Information Technology Allahabad,
India

The rapid integration of Large Language Models (LLMs) and
Retrieval-Augmented Generation (RAG) in academic research workflows has
introduced critical vulnerabilities regarding the provenance, accuracy,
and temporal validity of cited literature. This comprehensive survey
systematically evaluates automated compu-tational pipelines designed to
verify AI-generated references. We propose an extended,
multi-dimensional taxonomy that categorizes existing architectures into
metadata-matching heuristics, temporal graph-mining networks, and
multi-agent verification frameworks. By critically analyzing the
structural methodologies, mathematical formulations, and dataset
dependencies of these systems, we trace the paradigm shift from basic
lexical string-matching to deep semantic entailment via Natural Language
Inference (NLI). Furthermore, this review conducts a rigorous
comparative analysis of current frameworks against computational
bottlenecks, particularly API latency, dynamic graph update frequencies,
and paywalled knowledge enclosures. Finally, we outline a roadmap for
future research, emphasizing the development of localized Small Language
Model (SLM) infrastructures and heuristic models capable of detecting
softly superseded scientific methodologies.
\end{quote}

Additional Key Words and Phrases: Citation Verification, LLM
Hallucinations, Retrieval-Augmented Generation, Temporal Graph Networks,
Agentic Workflows, Semantic Fidelity, Natural Language Inference,
Scientometrics

\begin{quote}
\textbf{ACM Reference Format:}\\
Ritam Santra. 2026. An Automated Pipeline for Flagging Outdated
Citations in AI-Recommended Literature: A Comprehensive Review. 1, 1
(September 2026), 8 pages.

\textbf{1} \textbf{Introduction}

The landscape of scientific research is undergoing a fundamental
transformation driven by the advent of Generative Artificial
Intelligence (GenAI). Large Language Models (LLMs), such as OpenAI's
GPT-4, Anthropic's Claude, and open-weight models like LLaMA, serve as
powerful cognitive amplifiers. They enable researchers to rapidly
synthesize vast corpora of literature, draft academic manuscripts, and
discover cross-disciplinary connections that might otherwise remain
hidden. However, the foundational probabilistic architecture of these
models---which generate text by predicting statistically likely token
sequences rather than retrieving discrete, verified facts---inherently
leads to the generation of plausible but factually incorrect outputs
{[}\textbf{?} {]}. In academic contexts, this phenomenon manifests
dangerously as citation hallucinations.

Citation anomalies in AI-generated texts generally fall into three
distinct categories:

(1) \textbf{Complete Fabrication:} The authors, title, and venue are
entirely invented by the model, often blending recognizable names in the
field into a highly plausible but non-existent reference.

(2) \textbf{Entity Swapping:} Real authors are misattributed to real but
unrelated papers, or the DOI points to a completely different
manuscript.

Author's Contact Information: Ritam Santra, Indian Institute of
Information Technology Allahabad, Prayagraj, Uttar Pradesh, India.

Permission to make digital or hard copies of all or part of this work
for personal or classroom use is granted without fee provided that
copies are not made or distributed for profit or commercial advantage
and that copies bear this notice and the full citation on the first
page. Copyrights for components of this work owned by others than the
author(s) must be honored. Abstracting with credit is permitted. To copy
otherwise, or republish, to post on servers or to redistribute to lists,
requires prior specific permission and/or a fee. Request permissions
from permissions@acm.org.
\end{quote}

© 2026 Copyright held by the owner/author(s). Publication rights
licensed to ACM.

\begin{quote}
ACM XXXX-XXXX/2026/9-ART
\end{quote}

, Vol. 1, No. 1, Article . Publication date: September 2026.

\begin{quote}
2 Ritam Santra

(3) \textbf{Contextual Misalignment (Semantic Infidelity):} A genuine
paper is cited accurately with correct metadata, but the underlying text
of the paper does not support the specific claim generated by the LLM
{[}\textbf{?} {]}.

Even when LLMs are grounded using external knowledge bases via
Retrieval-Augmented Gener-ation (RAG) architectures, they frequently
retrieve and cite chronologically outdated or formally retracted
literature {[}\textbf{?} {]}. If a RAG system prioritizes dense vector
similarity over temporal metadata, it may cite a retracted study or a
methodology that has been universally superseded by the broader
scientific community.

The manual auditing of AI-generated bibliographies is a highly
resource-intensive bottleneck, often negating the efficiency gains
promised by AI assistance. Consequently, the development of robust,
automated computational pipelines to verify, contextualize, and flag
problematic citations has emerged as a critical imperative at the
intersection of Natural Language Processing (NLP) and computational
scientometrics.

\textbf{1.1} \textbf{Research Objectives and Contributions}
\end{quote}

This survey systematically analyzes the state-of-the-art in automated
citation verification. We investigate the core research question:
\emph{How can automated computational pipelines effectively verify,}

\emph{contextualize, and flag hallucinated, retracted, or
chronologically outdated citations generated by large}

\begin{quote}
\emph{language models?}

The primary contributions of this review are:

• A formalization of the citation hallucination problem, distinguishing
between parametric knowledge decay and semantic misalignment.

• A comprehensive, three-tier taxonomy of verification pipelines,
categorized into metadata, graph-based, and agentic approaches.

• A detailed analysis of the algorithmic frameworks powering these
systems, including Graph Neural Networks (GNNs) and Natural Language
Inference (NLI) models.\\
• A critical evaluation of current systemic bottlenecks, focusing on API
latency, metric limita-tions, and paywalled literature.

The remainder of this paper is structured as follows. Section 2
establishes the theoretical back-ground of LLM hallucinations and
formalizes the problem. Section 3 examines Tier 1 metadata and lexical
matching frameworks. Section 4 delves into Tier 2 topological and
temporal graph-mining approaches. Section 5 explores the frontier of
Tier 3 agentic verification and semantic fidelity. Section 6 outlines
evaluation metrics and benchmarking. Section 7 details open research
challenges, and Section 8 concludes the review.

\textbf{2} \textbf{Background and Problem Formalization}

\textbf{2.1} \textbf{The Anatomy of RAG Failures}
\end{quote}

To understand the necessity of automated verification pipelines, one
must examine the fundamental failure modes of Retrieval-Augmented
Generation (RAG). A standard RAG pipeline embeds a user query \emph{𝑞},
retrieves the top-\emph{𝑘}most similar document chunks \emph{𝐷}=
\{\emph{𝑑}1\emph{,𝑑}2\emph{, ...,𝑑𝑘}\} from a vector database using a
retriever function R, and prepends them to the LLM's context window to
generate a response \emph{𝑦}:\\
\emph{𝑦}= LLM(\emph{𝑞,} R(\emph{𝑞,𝑘})) (1)

However, semantic similarity in the dense embedding space does not
equate to scientific consen-sus or temporal validity {[}\textbf{?} {]}.
For example, if a query asks for "state-of-the-art image
classification," a standard RAG system might retrieve a highly cited
paper on Convolutional Neural Networks (CNNs) from 2015, ignoring a more
recent but less cited paper on Vision Transformers (ViTs). The

\begin{quote}
, Vol. 1, No. 1, Article . Publication date: September 2026.

An Automated Pipeline for Flagging Outdated Citations in AI-Recommended
Literature: A Comprehensive Review 3

LLM will then confidently cite the outdated 2015 paper. This is not a
hallucination of the text itself, but a hallucination of the
\emph{current scientific consensus}.

\textbf{2.2} \textbf{Formalizing Knowledge Decay and Semantic Drift}
\end{quote}

We can formalize the temporal validity of a cited fact. Let
K\emph{𝑡}represent the globally accepted scientific knowledge graph at
time \emph{𝑡}. A generated claim \emph{𝑐}based on a document
\emph{𝑑}published at time\emph{𝑡𝑑}is valid if and only if:

\begin{quote}
\emph{𝑐}\textbar= K\emph{𝑡𝑐𝑢𝑟𝑟}where \emph{𝑡𝑑}≤\emph{𝑡𝑐𝑢𝑟𝑟}(2)
\end{quote}

Automated pipelines must bridge the gap between \emph{𝑡𝑑}and
\emph{𝑡𝑐𝑢𝑟𝑟}. Semantic drift occurs when a claim \emph{𝑐}was valid in
K\emph{𝑡𝑑}but has been falsified or superseded in K\emph{𝑡𝑐𝑢𝑟𝑟}.

\begin{quote}
\textbf{2.3} \textbf{A Comprehensive Taxonomy}
\end{quote}

We propose an extended taxonomy that categorizes automated verification
pipelines into three primary architectural tiers, moving from structural
validation to deep semantic reasoning:

\begin{quote}
(1) \textbf{Tier 1: Metadata-Centric (Lexical) Approaches:} Systems
relying on exact or fuzzy match- ing of citation strings against
established indexing APIs (e.g., Crossref, OpenAlex).

(2) \textbf{Tier 2: Topological and Graph-Mining Approaches:} Systems
that map the citation network as a Directed Acyclic Graph (DAG) to
verify logical existence, temporal relevance, and community centrality.

(3) \textbf{Tier 3: Agentic and LLM-in-the-Loop Approaches:} Systems
that deploy secondary LLM agents to semantically read and verify the
source text against the generated claim using NLI.

\textbf{3} \textbf{Tier 1: Metadata and Lexical Matching Frameworks}

The foundational layer of automated verification relies on metadata
matching. These pipelines extract the raw citation string from the LLM
output and attempt to resolve it against external databases.

\textbf{3.1} \textbf{Citation Parsing and Extraction}

Before verification can occur, the pipeline must accurately segment the
generated text. Tools like GROBID (GeneRation Of BIbliographic Data)
utilize Conditional Random Fields (CRF) and, more recently, deep
learning sequence-to-sequence models to parse raw citation strings into
structured JSON objects containing fields like \emph{author},
\emph{title}, \emph{journal}, and \emph{year}.

\textbf{3.2} \textbf{String Similarity and Fuzzy Matching}
\end{quote}

Wang et al. {[}\textbf{?} {]} formalized the metadata matching approach
in \emph{RefChecker}. Once parsed, the pipeline queries an API (like
Crossref) using the title. The core of lexical matching relies on
calculating string similarities between the generated citation
\emph{𝐶𝑔𝑒𝑛}and the retrieved database record \emph{𝐶𝑑𝑏}. Common metrics
include Jaccard similarity:

\begin{longtable}[]{@{}
  >{\raggedright\arraybackslash}p{(\columnwidth - 0\tabcolsep) * \real{1.0000}}@{}}
\toprule()
\begin{minipage}[b]{\linewidth}\raggedright
\begin{quote}
\emph{𝐽}(\emph{𝐶𝑔𝑒𝑛,𝐶𝑑𝑏}) =\textbar{}\emph{𝐶𝑔𝑒𝑛}∩\emph{𝐶𝑑𝑏}\textbar{}
\textbar{}\emph{𝐶𝑔𝑒𝑛}∪\emph{𝐶𝑑𝑏}\textbar{} (3)

To handle character-level noise and hallucinated middle initials
generated by LLMs, the Leven-

shtein distance \emph{𝑙𝑒𝑣}(\emph{𝑎,𝑏}) is often normalized into a
similarity score:

\emph{𝑆𝑖𝑚𝑙𝑒𝑣}(\emph{𝐶𝑔𝑒𝑛,𝐶𝑑𝑏}) = 1
−max(\textbar{}\emph{𝐶𝑔𝑒𝑛}\textbar{}\emph{,}
\textbar{}\emph{𝐶𝑑𝑏}\textbar) \emph{𝑙𝑒𝑣}(\emph{𝐶𝑔𝑒𝑛,𝐶𝑑𝑏}) (4)
\end{quote}

If \emph{𝑆𝑖𝑚𝑙𝑒𝑣}≥\emph{𝜏}(where \emph{𝜏}is a strict threshold, typically
0\emph{.}85), the citation is deemed structurally valid.
\end{minipage} \\
\midrule()
\endhead
\bottomrule()
\end{longtable}

, Vol. 1, No. 1, Article . Publication date: September 2026.

\begin{longtable}[]{@{}
  >{\raggedright\arraybackslash}p{(\columnwidth - 2\tabcolsep) * \real{0.5000}}
  >{\raggedright\arraybackslash}p{(\columnwidth - 2\tabcolsep) * \real{0.5000}}@{}}
\toprule()
\begin{minipage}[b]{\linewidth}\raggedright
\begin{quote}
4
\end{quote}
\end{minipage} & \begin{minipage}[b]{\linewidth}\raggedright
Ritam Santra
\end{minipage} \\
\midrule()
\endhead
\bottomrule()
\end{longtable}

\begin{longtable}[]{@{}
  >{\raggedright\arraybackslash}p{(\columnwidth - 0\tabcolsep) * \real{1.0000}}@{}}
\toprule()
\begin{minipage}[b]{\linewidth}\raggedright
\end{minipage} \\
\midrule()
\endhead
\bottomrule()
\end{longtable}

Fig. 1. Comprehensive Taxonomy and Workflow Architecture of Tier 1, Tier
2, and Tier 3 Citation Verification

Pipelines. The diagram illustrates the data flow from claim generation
to metadata parsing, graph traversal,

\begin{quote}
and final semantic entailment.

\textbf{3.3} \textbf{Limitations of Tier 1}
\end{quote}

While computationally inexpensive (requiring only milliseconds per
citation) and highly scalable, metadata matching is inherently brittle.
It successfully identifies completely fabricated citations but fails
catastrophically at identifying Contextual Misalignment. A Tier 1 system
will verify a citation as perfectly accurate as long as the paper exists
in the database, regardless of whether the paper actually supports the
LLM's claim.

\begin{quote}
\textbf{4} \textbf{Tier 2: Topological and Temporal Graph Networks}

To move beyond structural existence, Tier 2 frameworks model academic
literature as complex networks. Let G = (\emph{𝑉, 𝐸}) be a directed
citation graph where \emph{𝑉}is the set of papers and \emph{𝐸}is the set
of directed edges (citations).

\textbf{4.1} \textbf{Detecting Anomalous Citations via GNNs}

Zhang and Shi {[}\textbf{?} {]} demonstrated that hallucinated citations
exhibit highly anomalous topological features. Real citations form
localized, dense clusters within specific sub-disciplines. If an LLM
hallucinates a connection between a paper in quantum mechanics and a
paper in medieval history, the shortest path \emph{𝑑}(\emph{𝑢, 𝑣}) in
the citation graph will be abnormally large.
\end{quote}

edge existing based on structural node similarity. Using Graph
Convolutional Networks (GCN), the Tier 2 pipelines utilize Graph Neural
Networks (GNNs) to calculate the probability \emph{𝑃}(\emph{𝑒𝑢,𝑣}) of an

\begin{quote}
, Vol. 1, No. 1, Article . Publication date: September 2026.

An Automated Pipeline for Flagging Outdated Citations in AI-Recommended
Literature: A Comprehensive Review 5

node embeddings H(\emph{𝑙})at layer \emph{𝑙}are updated as:
\end{quote}

where˜A is the adjacency matrix with added
self-connections,˜H(\emph{𝑙}+1)= \emph{𝜎}(︂D−1 2 ˜A˜D−1 2
H(\emph{𝑙})W(\emph{𝑙}))︂D is the degree matrix, and W is the (5)

\begin{quote}
learned weight matrix. Citations generated by the LLM that result in a
predicted edge probability\emph{𝑃}(\emph{𝑒𝑢,𝑣}) \emph{\textless{} 𝜖}are
flagged as likely hallucinations for human review.

\textbf{4.2} \textbf{Temporal Decay and Outdated Literature}

Flagging outdated citations requires incorporating continuous temporal
dynamics into the graph. Kim and Park {[}\textbf{?} {]} introduced the
LAGMiD framework, which utilizes a temporal decay function. The ongoing
relevance \emph{𝑅}of a cited paper \emph{𝑣}at the current evaluation
time \emph{𝑡𝑐𝑢𝑟𝑟}is formulated as:

\emph{𝑅}(\emph{𝑣,𝑡𝑐𝑢𝑟𝑟}) = PageRank(\emph{𝑣,} G) · exp
(−\emph{𝜆}(\emph{𝑡𝑐𝑢𝑟𝑟}−\emph{𝑡𝑣})) (6)

where\emph{𝑡𝑣}is the publication year of the paper and \emph{𝜆}is a
domain-specific decay constant (e.g., computer science literature decays
faster than mathematics literature).

\textbf{Algorithm 1} LAGMiD Temporal Graph Verification Algorithm
\end{quote}

\begin{longtable}[]{@{}
  >{\raggedright\arraybackslash}p{(\columnwidth - 2\tabcolsep) * \real{0.5000}}
  >{\raggedright\arraybackslash}p{(\columnwidth - 2\tabcolsep) * \real{0.5000}}@{}}
\toprule()
\multicolumn{2}{@{}>{\raggedright\arraybackslash}p{(\columnwidth - 2\tabcolsep) * \real{1.0000} + 2\tabcolsep}@{}}{%
\begin{minipage}[b]{\linewidth}\raggedright
\textbf{Require:} Set of generated citations \emph{𝐶𝑔𝑒𝑛}, Citation Graph
G = (\emph{𝑉, 𝐸}), Current Year \emph{𝑡𝑐𝑢𝑟𝑟}, Decay threshold \emph{𝜃}

\begin{quote}
\textbf{Ensure:} Set of flagged outdated citations \emph{𝐹𝑜𝑢𝑡𝑑𝑎𝑡𝑒𝑑}

1: \emph{𝐹𝑜𝑢𝑡𝑑𝑎𝑡𝑒𝑑}←∅\\
2: \textbf{for} each \emph{𝑐}∈\emph{𝐶𝑔𝑒𝑛}\textbf{do} 3:

4:\emph{𝑣}←MatchNode(\emph{𝑐,} G) \textbf{if} \emph{𝑣}== NULL
\textbf{then}

5: Flag \emph{𝑐}as Fabricated
\end{quote}\strut
\end{minipage}} \\
\midrule()
\endhead
6: 7: 8: 9: 10: & \begin{minipage}[t]{\linewidth}\raggedright
\begin{quote}
\textbf{end if}

\emph{𝑃𝑅𝑣}←CalculatePageRank(\emph{𝑣,} G)\\
\emph{𝑡𝑣}←GetPublicationYear(\emph{𝑣})\\
\emph{𝑅𝑣}←\emph{𝑃𝑅𝑣}· exp(−\emph{𝜆}(\emph{𝑡𝑐𝑢𝑟𝑟}−\emph{𝑡𝑣})) \textbf{if}
\emph{𝑅𝑣\textless{} 𝜃}\textbf{then}
\end{quote}\strut
\end{minipage} \\
\multicolumn{2}{@{}>{\raggedright\arraybackslash}p{(\columnwidth - 2\tabcolsep) * \real{1.0000} + 2\tabcolsep}@{}}{%
\begin{minipage}[t]{\linewidth}\raggedright
\begin{quote}
11: 12: \textbf{end if}\emph{𝐹𝑜𝑢𝑡𝑑𝑎𝑡𝑒𝑑}←\emph{𝐹𝑜𝑢𝑡𝑑𝑎𝑡𝑒𝑑}∪\{\emph{𝑐}\}

13: \textbf{end for}\\
14: \textbf{return} \emph{𝐹𝑜𝑢𝑡𝑑𝑎𝑡𝑒𝑑}
\end{quote}\strut
\end{minipage}} \\
\begin{minipage}[t]{\linewidth}\raggedright
\begin{quote}
\textbf{5}
\end{quote}
\end{minipage} & \begin{minipage}[t]{\linewidth}\raggedright
\begin{quote}
\textbf{Tier 3: Agentic Verification and Semantic Fidelity}
\end{quote}
\end{minipage} \\
\bottomrule()
\end{longtable}

\begin{quote}
The vanguard of citation verification relies on multi-agent
architectures that simulate human peer-review. Tier 3 systems fetch the
source document and computationally "read" it to ensure it semantically
entails the generated claim.

\textbf{5.1} \textbf{Natural Language Inference (NLI) Modeling}

Gupta et al. {[}\textbf{?} {]} established the Semantic Fidelity (SemFi)
pipeline, framing verification as an NLI classification task. Given an
LLM-generated claim \emph{𝑐}and the retrieved abstract or full-text of
the cited document \emph{𝑑}, a cross-encoder model (e.g., RoBERTa or
DeBERTa) computes the probability of entailment.
\end{quote}

, Vol. 1, No. 1, Article . Publication date: September 2026.

\begin{quote}
6 Ritam Santra

The input is formatted as `{[}CLS{]} claim {[}SEP{]} document
{[}SEP{]}`, and the pooled output is passed through a Multi-Layer
Perceptron (MLP):

\emph{𝑃}(Class \textbar{} \emph{𝑐,𝑑}) = Softmax(MLP(E\emph{𝐶𝐿𝑆})) (7)
\end{quote}

The output classifies the relationship as \emph{Entailment} (valid
citation, the document supports the claim), \emph{Neutral} (tangential
citation, the document mentions the topic but does not prove the claim),
or \emph{Contradiction} (hallucinated citation, the document refutes the
claim).

\begin{quote}
\textbf{5.2} \textbf{Iterative Reflection (ReAct) Workflows}

Li et al. {[}\textbf{?} {]} expanded NLI into an autonomous agentic
framework using the ReAct (Reasoning and Acting) prompting paradigm
{[}\textbf{?} {]}. In the BibAgent system, if the agent determines the
cited paper does not support the claim (a \emph{Contradiction} or
\emph{Neutral} score), it does not merely flag the error. Instead, it
enters an autonomous reflection loop.

The agent formulates a new search query based on the unsupported claim,
queries the Semantic Scholar API, downloads candidate abstracts, and
evaluates them sequentially until a valid paper is found that
\emph{does} entail the claim. It then silently replaces the LLM's
hallucinated citation with the correct, verified citation before
presenting the final manuscript to the user.
\end{quote}

\begin{longtable}[]{@{}
  >{\raggedright\arraybackslash}p{(\columnwidth - 0\tabcolsep) * \real{1.0000}}@{}}
\toprule()
\begin{minipage}[b]{\linewidth}\raggedright
\end{minipage} \\
\midrule()
\endhead
\bottomrule()
\end{longtable}

Fig. 2. The iterative ReAct loop of an Agentic Verification system. The
agent oscillates between reasoning

\begin{quote}
about entailment and acting upon search APIs until a valid citation is
secured.

\textbf{6} \textbf{Evaluation Metrics and Benchmarking}

Evaluating the efficacy of automated verification pipelines requires
specialized metrics beyond standard NLP benchmarks. Current literature
relies on a combination of classification accuracy and operational
efficiency.

\textbf{6.1} \textbf{Classification Metrics}

For detecting fabricated and outdated citations, systems are evaluated
on standard Precision, Recall, and F1 scores against manually annotated
datasets (such as the HaluEval-Sci benchmark).
\end{quote}

\begin{longtable}[]{@{}
  >{\raggedright\arraybackslash}p{(\columnwidth - 10\tabcolsep) * \real{0.1667}}
  >{\raggedright\arraybackslash}p{(\columnwidth - 10\tabcolsep) * \real{0.1667}}
  >{\raggedright\arraybackslash}p{(\columnwidth - 10\tabcolsep) * \real{0.1667}}
  >{\raggedright\arraybackslash}p{(\columnwidth - 10\tabcolsep) * \real{0.1667}}
  >{\raggedright\arraybackslash}p{(\columnwidth - 10\tabcolsep) * \real{0.1667}}
  >{\raggedright\arraybackslash}p{(\columnwidth - 10\tabcolsep) * \real{0.1667}}@{}}
\toprule()
\begin{minipage}[b]{\linewidth}\raggedright
Precision =
\end{minipage} & \begin{minipage}[b]{\linewidth}\raggedright
\begin{quote}
\emph{𝑇𝑃}\\
\emph{𝑇𝑃}+ \emph{𝐹𝑃,}
\end{quote}\strut
\end{minipage} & \begin{minipage}[b]{\linewidth}\raggedright
\begin{quote}
Recall =
\end{quote}
\end{minipage} & \begin{minipage}[b]{\linewidth}\raggedright
\begin{quote}
\emph{𝑇𝑃}\\
\emph{𝑇𝑃}+ \emph{𝐹𝑁,}
\end{quote}\strut
\end{minipage} & \begin{minipage}[b]{\linewidth}\raggedright
\begin{quote}
\emph{𝐹}1 = 2 ·Precision · Recall Precision + Recall
\end{quote}
\end{minipage} & \begin{minipage}[b]{\linewidth}\raggedright
(8)
\end{minipage} \\
\midrule()
\endhead
\bottomrule()
\end{longtable}

\begin{quote}
, Vol. 1, No. 1, Article . Publication date: September 2026.

An Automated Pipeline for Flagging Outdated Citations in AI-Recommended
Literature: A Comprehensive Review 7
\end{quote}

Where True Positives (\emph{𝑇𝑃}) represent successfully flagged
hallucinations or outdated references.

\begin{quote}
\textbf{6.2} \textbf{Operational Metrics (Latency)}

A critical metric for Tier 3 systems is End-to-End Latency
(L\emph{𝑒}2\emph{𝑒}). If an LLM generates a document with
\emph{𝑁}citations, the total verification time is bounded by the network
latency of the API calls (\emph{𝑡𝑎𝑝𝑖}) and the inference time of the NLI
model (\emph{𝑡𝑖𝑛𝑓}):

where \emph{𝑘}is the number of iterative reflection loops required per
citation. L\emph{𝑒}2\emph{𝑒}=∑︁(︁\emph{𝑘}· \emph{𝑡𝑎𝑝𝑖}+
\emph{𝑡𝑖𝑛𝑓}(\textbar{}\emph{𝑑𝑖}\textbar))︁ (9)
\end{quote}

Table 1. Comprehensive Evaluation and Comparative Analysis of Automated
Citation Verification Pipelines

\begin{longtable}[]{@{}
  >{\raggedright\arraybackslash}p{(\columnwidth - 20\tabcolsep) * \real{0.0909}}
  >{\raggedright\arraybackslash}p{(\columnwidth - 20\tabcolsep) * \real{0.0909}}
  >{\raggedright\arraybackslash}p{(\columnwidth - 20\tabcolsep) * \real{0.0909}}
  >{\raggedright\arraybackslash}p{(\columnwidth - 20\tabcolsep) * \real{0.0909}}
  >{\raggedright\arraybackslash}p{(\columnwidth - 20\tabcolsep) * \real{0.0909}}
  >{\raggedright\arraybackslash}p{(\columnwidth - 20\tabcolsep) * \real{0.0909}}
  >{\raggedright\arraybackslash}p{(\columnwidth - 20\tabcolsep) * \real{0.0909}}
  >{\raggedright\arraybackslash}p{(\columnwidth - 20\tabcolsep) * \real{0.0909}}
  >{\raggedright\arraybackslash}p{(\columnwidth - 20\tabcolsep) * \real{0.0909}}
  >{\raggedright\arraybackslash}p{(\columnwidth - 20\tabcolsep) * \real{0.0909}}
  >{\raggedright\arraybackslash}p{(\columnwidth - 20\tabcolsep) * \real{0.0909}}@{}}
\toprule()
\multicolumn{2}{@{}>{\raggedright\arraybackslash}p{(\columnwidth - 20\tabcolsep) * \real{0.1818} + 2\tabcolsep}}{%
\begin{minipage}[b]{\linewidth}\raggedright
\textbf{Framework}
\end{minipage}} &
\multicolumn{3}{>{\raggedright\arraybackslash}p{(\columnwidth - 20\tabcolsep) * \real{0.2727} + 4\tabcolsep}}{%
\begin{minipage}[b]{\linewidth}\raggedright
\begin{quote}
\textbf{Architecture}
\end{quote}
\end{minipage}} &
\multicolumn{2}{>{\raggedright\arraybackslash}p{(\columnwidth - 20\tabcolsep) * \real{0.1818} + 2\tabcolsep}}{%
\begin{minipage}[b]{\linewidth}\raggedright
\textbf{Target API}
\end{minipage}} &
\multirow{2}{*}{\begin{minipage}[b]{\linewidth}\raggedright
\begin{quote}
\textbf{Latency (}L\textbf{\_}\emph{𝑒}2\emph{𝑒}\textbf{)}\\
Very\\
(\emph{\textless{}}1s)
\end{quote}\strut
\end{minipage}} & \begin{minipage}[b]{\linewidth}\raggedright
\end{minipage} & \begin{minipage}[b]{\linewidth}\raggedright
\begin{quote}
\textbf{Key Capabilities}
\end{quote}
\end{minipage} & \begin{minipage}[b]{\linewidth}\raggedright
\begin{quote}
\textbf{Critical Limitations}
\end{quote}
\end{minipage} \\
\multicolumn{2}{@{}>{\raggedright\arraybackslash}p{(\columnwidth - 20\tabcolsep) * \real{0.1818} + 2\tabcolsep}}{%
\begin{minipage}[b]{\linewidth}\raggedright
\begin{quote}
RefChecker {[}\textbf{?} {]}
\end{quote}
\end{minipage}} &
\multicolumn{3}{>{\raggedright\arraybackslash}p{(\columnwidth - 20\tabcolsep) * \real{0.2727} + 4\tabcolsep}}{%
\begin{minipage}[b]{\linewidth}\raggedright
\begin{quote}
Tier 1 (Lexical)
\end{quote}
\end{minipage}} &
\multicolumn{2}{>{\raggedright\arraybackslash}p{(\columnwidth - 20\tabcolsep) * \real{0.1818} + 2\tabcolsep}}{%
\begin{minipage}[b]{\linewidth}\raggedright
\begin{quote}
Crossref
\end{quote}
\end{minipage}} & & \begin{minipage}[b]{\linewidth}\raggedright
Low
\end{minipage} & \begin{minipage}[b]{\linewidth}\raggedright
\begin{quote}
Highly scalable; reliably detects complete fabrica-tions.
\end{quote}
\end{minipage} & \begin{minipage}[b]{\linewidth}\raggedright
\begin{quote}
Fails entirely on seman-tic contextual errors; high false positive rate
on preprints.
\end{quote}
\end{minipage} \\
\midrule()
\endhead
\multicolumn{2}{@{}>{\raggedright\arraybackslash}p{(\columnwidth - 20\tabcolsep) * \real{0.1818} + 2\tabcolsep}}{%
LAGMiD {[}\textbf{?} {]}} &
\multicolumn{3}{>{\raggedright\arraybackslash}p{(\columnwidth - 20\tabcolsep) * \real{0.2727} + 4\tabcolsep}}{%
\begin{minipage}[t]{\linewidth}\raggedright
\begin{quote}
Tier 2 (Topologi-cal)
\end{quote}
\end{minipage}} &
\multicolumn{2}{>{\raggedright\arraybackslash}p{(\columnwidth - 20\tabcolsep) * \real{0.1818} + 2\tabcolsep}}{%
\begin{minipage}[t]{\linewidth}\raggedright
\begin{quote}
S2AG Graphs
\end{quote}
\end{minipage}} &
\multicolumn{2}{>{\raggedright\arraybackslash}p{(\columnwidth - 20\tabcolsep) * \real{0.1818} + 2\tabcolsep}}{%
\begin{minipage}[t]{\linewidth}\raggedright
\begin{quote}
Medium\\
(2-5s)
\end{quote}\strut
\end{minipage}} & \begin{minipage}[t]{\linewidth}\raggedright
\begin{quote}
Effectively maps chronologies and flags outdated work using decay
functions.
\end{quote}
\end{minipage} & \begin{minipage}[t]{\linewidth}\raggedright
\begin{quote}
Computationally expen-sive to update the dy-namic graph in real-time.
\end{quote}
\end{minipage} \\
\multicolumn{2}{@{}>{\raggedright\arraybackslash}p{(\columnwidth - 20\tabcolsep) * \real{0.1818} + 2\tabcolsep}}{%
\begin{minipage}[t]{\linewidth}\raggedright
\begin{quote}
sciwrite-lint {[}\textbf{?} {]}
\end{quote}
\end{minipage}} & \begin{minipage}[t]{\linewidth}\raggedright
Tier\\
RAG)\strut
\end{minipage} & 2 & \begin{minipage}[t]{\linewidth}\raggedright
\begin{quote}
(Local
\end{quote}
\end{minipage} & \begin{minipage}[t]{\linewidth}\raggedright
\begin{quote}
PubMed ArXiv
\end{quote}
\end{minipage} & / &
\multicolumn{2}{>{\raggedright\arraybackslash}p{(\columnwidth - 20\tabcolsep) * \real{0.1818} + 2\tabcolsep}}{%
\begin{minipage}[t]{\linewidth}\raggedright
\begin{quote}
Low (1-2s)
\end{quote}
\end{minipage}} & \begin{minipage}[t]{\linewidth}\raggedright
\begin{quote}
Ensures data privacy; high throughput for local checking without API
rate limits.
\end{quote}
\end{minipage} & \begin{minipage}[t]{\linewidth}\raggedright
\begin{quote}
Constrained by local

hardware memory;

limited disciplinary

scope.
\end{quote}
\end{minipage} \\
\multicolumn{2}{@{}>{\raggedright\arraybackslash}p{(\columnwidth - 20\tabcolsep) * \real{0.1818} + 2\tabcolsep}}{%
BibAgent {[}\textbf{?} {]}} &
\multicolumn{3}{>{\raggedright\arraybackslash}p{(\columnwidth - 20\tabcolsep) * \real{0.2727} + 4\tabcolsep}}{%
Tier 3 (Agentic)} &
\multicolumn{2}{>{\raggedright\arraybackslash}p{(\columnwidth - 20\tabcolsep) * \real{0.1818} + 2\tabcolsep}}{%
\begin{minipage}[t]{\linewidth}\raggedright
\begin{quote}
OpenAlex
\end{quote}
\end{minipage}} & \begin{minipage}[t]{\linewidth}\raggedright
\begin{quote}
Very\\
(\emph{\textgreater{}}30s)
\end{quote}\strut
\end{minipage} & High & \begin{minipage}[t]{\linewidth}\raggedright
\begin{quote}
High accuracy in deep se-mantic traceability; self-corrects
hallucinations automatically.
\end{quote}
\end{minipage} & \begin{minipage}[t]{\linewidth}\raggedright
\begin{quote}
Severe API latency dur-ing the agent's iterative reflection loops breaks
real-time UX.
\end{quote}
\end{minipage} \\
\multicolumn{2}{@{}>{\raggedright\arraybackslash}p{(\columnwidth - 20\tabcolsep) * \real{0.1818} + 2\tabcolsep}}{%
\begin{minipage}[t]{\linewidth}\raggedright
\begin{quote}
RetractDet {[}\textbf{?} {]}
\end{quote}
\end{minipage}} &
\multicolumn{3}{>{\raggedright\arraybackslash}p{(\columnwidth - 20\tabcolsep) * \real{0.2727} + 4\tabcolsep}}{%
\begin{minipage}[t]{\linewidth}\raggedright
\begin{quote}
Tier 1 (Filter)
\end{quote}
\end{minipage}} &
\multicolumn{2}{>{\raggedright\arraybackslash}p{(\columnwidth - 20\tabcolsep) * \real{0.1818} + 2\tabcolsep}}{%
\begin{minipage}[t]{\linewidth}\raggedright
\begin{quote}
Retraction Watch
\end{quote}
\end{minipage}} &
\multicolumn{2}{>{\raggedright\arraybackslash}p{(\columnwidth - 20\tabcolsep) * \real{0.1818} + 2\tabcolsep}}{%
\begin{minipage}[t]{\linewidth}\raggedright
\begin{quote}
Low (\emph{\textless{}}1s)
\end{quote}
\end{minipage}} & \begin{minipage}[t]{\linewidth}\raggedright
\begin{quote}
Instantly eliminates for-mally retracted and dan-gerous sources from
gen-eration.
\end{quote}
\end{minipage} & \begin{minipage}[t]{\linewidth}\raggedright
\begin{quote}
Highly dependent on ex-ternal database update frequencies and manual
flagging.
\end{quote}
\end{minipage} \\
\multicolumn{2}{@{}>{\raggedright\arraybackslash}p{(\columnwidth - 20\tabcolsep) * \real{0.1818} + 2\tabcolsep}}{%
\begin{minipage}[t]{\linewidth}\raggedright
\begin{quote}
SemFi {[}\textbf{?} {]}
\end{quote}
\end{minipage}} &
\multicolumn{3}{>{\raggedright\arraybackslash}p{(\columnwidth - 20\tabcolsep) * \real{0.2727} + 4\tabcolsep}}{%
\begin{minipage}[t]{\linewidth}\raggedright
\begin{quote}
Tier 3 (NLI)
\end{quote}
\end{minipage}} &
\multicolumn{2}{>{\raggedright\arraybackslash}p{(\columnwidth - 20\tabcolsep) * \real{0.1818} + 2\tabcolsep}}{%
\begin{minipage}[t]{\linewidth}\raggedright
\begin{quote}
Custom Cor-pus
\end{quote}
\end{minipage}} &
\multicolumn{2}{>{\raggedright\arraybackslash}p{(\columnwidth - 20\tabcolsep) * \real{0.1818} + 2\tabcolsep}}{%
High (10-15s)} & \begin{minipage}[t]{\linewidth}\raggedright
\begin{quote}
Solves contextual hallu-cination via robust cross-encoder entailment
mod-eling.
\end{quote}
\end{minipage} & \begin{minipage}[t]{\linewidth}\raggedright
\begin{quote}
Struggles with highly complex mathematical equations or domain-specific
jargon in the full text.
\end{quote}
\end{minipage} \\
\textbf{7} &
\multicolumn{8}{>{\raggedright\arraybackslash}p{(\columnwidth - 20\tabcolsep) * \real{0.7273} + 14\tabcolsep}}{%
\begin{minipage}[t]{\linewidth}\raggedright
\begin{quote}
\textbf{Research Gaps and Open Challenges}
\end{quote}
\end{minipage}} & & \\
\bottomrule()
\end{longtable}

\begin{quote}
Despite rapid architectural innovations, several critical bottlenecks
prevent these pipelines from achieving seamless integration into
everyday academic drafting workflows.

\textbf{7.1} \textbf{The API Latency and Rate-Limit Bottleneck}

As identified by O'Connor {[}\textbf{?} {]}, real-time verification is
currently paralyzed by API rate limits. When a generative model produces
a comprehensive literature review containing 50 to 100 citations, a Tier
3 agentic framework like BibAgent must execute hundreds of sequential
API calls to fetch metadata, download full-text abstracts, embed the
documents, and run NLI inference. This cascading
\end{quote}

, Vol. 1, No. 1, Article . Publication date: September 2026.

\begin{quote}
8 Ritam Santra

requirement results in processing times that can exceed ten minutes per
document, effectively breaking the real-time, interactive user
experience researchers expect from AI tools.

\textbf{7.2} \textbf{The Nuance of "Softly" Outdated Literature}

Current pipelines excel at identifying binary states: a paper is either
real or fabricated, formally retracted or active. However, academic
obsolescence is rarely binary; it exists on a continuous gradient. A
paper from 2017 detailing an NLP methodology using LSTMs is not
retracted, but it has been scientifically superseded by
Transformer-based architectures. Current graph-mining approaches (such
as Equation 6) rely on rigid mathematical decay, failing to capture the
semantic nuance of paradigm shifts. Developing heuristic models that can
read the context of a citation (e.g., citing an LSTM paper for
historical context vs. citing it as the current state-of-the-art
methodology) remains a significant open challenge.

\textbf{7.3} \textbf{Paywalled Knowledge Enclosures}
\end{quote}

Tier 3 agentic verification strictly requires access to the full text of
a paper to verify deep semantic fidelity. However, a significant portion
of high-impact academic literature remains locked behind publisher
paywalls. Consequently, verification pipelines are forced to fall back
on analyzing publicly available abstracts. Because abstracts often omit
granular methodological details, dataset specifics, and nuanced
limitations, agentic models frequently trigger false negatives, assuming
a citation is hallucinated simply because the supporting evidence was
buried in the inaccessible full-text methodology section.

\begin{quote}
\textbf{7.4} \textbf{Multimodal and Data-Driven Citations}

Current verification frameworks are entirely text-centric. However,
modern scientific discourse increasingly relies on multimodal citations,
where claims are based on specific figures, data tables, or open-source
repositories linked within a paper. An LLM might correctly cite a paper
but misinterpret a chart within that paper, generating a false claim.
Extending NLI models to parse multimodal PDFs and verify data-centric
claims represents a largely unexplored frontier in computational
scientometrics.

\textbf{8} \textbf{Conclusion}
\end{quote}

The automated verification of AI-generated citations is an urgent
imperative for preserving the integrity of the scientific record in the
era of Generative AI. This survey has charted the rapid evolution of
verification pipelines from basic metadata string-matching (Tier 1) to
topological graph mapping (Tier 2) and sophisticated, LLM-driven
semantic entailment (Tier 3).

\begin{quote}
While tools like RefChecker established a scalable baseline, they are
insufficient for catching con-textual hallucinations. Conversely,
frameworks like BibAgent and SemFi offer remarkable accuracy but are
severely hampered by API latency and paywalled text restrictions. Future
research must prioritize bridging the gap between deep semantic
verification and computational efficiency. The development of localized,
lightweight Small Language Models (SLMs) that operate on compressed,
open-access vector indices will be crucial. Furthermore, integrating
context-aware temporal dy-namics---allowing pipelines to intelligently
differentiate between foundational historical citations and superseded
obsolete methodologies---will define the next generation of automated
academic auditing.

\textbf{References}

, Vol. 1, No. 1, Article . Publication date: September 2026.
\end{quote}

\end{document}
